% Conclusión de: Gilces Carranza José Ignacio
\subsection{Visión y alcance validados con la organización}

\textbf{Afirmación.} La visión y el alcance de un sistema solo son confiables después de contrastarlos con la organización real, aunque el equipo ya los tenga redactados.

\textbf{Evidencia.} La visión de la práctica de la Unidad~I prometía sugerir sustitutos de forma automática. En las entrevistas del 19 de noviembre la cajera con experiencia y el propietario rechazaron esa función por el riesgo de dañar el equipo del cliente (Sección~\ref{sec:ficha}). La norma ISO/IEC/IEEE 29148:2018 pide definir el límite del sistema a partir de las partes interesadas antes de especificar requisitos~\cite{iso29148}, y seguir ese orden permitió detectar el problema a tiempo.

\textbf{Implicación.} La PVIA informa y sugiere complementos, pero la decisión de sustituir queda en el vendedor (RE-02). La emisión de comprobantes ante el SRI queda en zona gris hasta la reunión con el propietario de la semana~6; ningún requisito sobre ese tema se especificará antes.
